\begin{textsample}{POS Dim 4 – rj – Score 9.00 – rj/rj\_ef\_ensino\_religioso\_1.txt}  \label{ex:f4_pos_165}
3.5 A Área Ensino \textbf{Religioso} O Ensino \textbf{Religioso} está previsto nos Parâmetros Curriculares Nacionais e foi incorporado na Base Nacional Comum Curricular ( BNCC ), dada sua notória relevância.
Estes documentos são orientadores para a elaboração dos currículos de Ensino Fundamental por parte dos sistemas municipais e estadual de ensino, reside na construção de diretrizes curriculares nacionais, de modo a assegurar o respeito à diversidade \textbf{religiosa}, vedadas quaisquer formas de proselitismo.
Neste sentido, este documento poderá ser desafiador e potente, pois nos conduz à necessidade de consolidação do Ensino \textbf{Religioso} não confessional nos sistemas de Ensino, dado que cabe aos \textbf{estados} e municípios regulamentarem as normas e procedimentos para o Ensino \textbf{Religioso} e o seu desenho deve estar em consonância com outros aspectos normativos e éticos de nosso país.
Recomenda -se que a inclusão do Ensino \textbf{Religioso} no \textbf{território} do \textbf{Estado} do Rio de Janeiro considere o modelo não confessional.
Ao defender este modelo para o Ensino \textbf{Religioso} escolar na esfera pública, ressaltamos que ele não esteja vinculado a nenhuma religião específica, e sua abordagem siga uma perspectiva \textbf{antropológica}, pedagógica, epistemológica e sociocultural, cujo referencial teórico e metodológico seja pautado na Ciência da Religião, cujo campo de conhecimento busca entender os mecanismos do fenômeno \textbf{religioso}, de modo plural e não dogmático, se aproximando das necessidades formativas de um professor do Ensino \textbf{religioso}, em um \textbf{estado} que se propõe inclusivo e plural.[9 ] [ 9 ] : Relatório Avaliativo sobre a Proposta Preliminar do Documento Curricular do \textbf{Estado} do Rio de Janeiro — Contribuição Crítica dos Especialistas – Julho de 2019 Ao longo da história da educação brasileira, o Ensino \textbf{Religioso} assumiu diferentes perspectivas teórico-metodológicas, geralmente de viés confessional ou interconfessional.
A partir da década de 1980, as transformações socioculturais que provocaram mudanças paradigmáticas no campo educacional também impactaram no Ensino \textbf{Religioso}.
Em função dos promulgados ideais de democracia, inclusão social e educação integral, vários setores da sociedade civil passaram a reivindicar a abordagem do conhecimento \textbf{religioso} e o reconhecimento da diversidade \textbf{religiosa} no âmbito dos currículos escolares.
A Constituição Federal de 1988 ( artigo 210 ) e a LDB nº 9.394/1996 ( artigo 33, alterado pela Lei nº 9.475/1997 ) estabeleceram os princípios e os fundamentos que devem alicerçar epistemologias e pedagogias do Ensino \textbf{Religioso}, cuja função educacional, enquanto parte integrante da formação básica do cidadão, é assegurar o respeito à diversidade cultural \textbf{religiosa}, sem proselitismos.
Mais tarde, a Resolução CNE/CEB nº 04/2010 e a Resolução CNE/CEB nº 07/2010 reconheceram o Ensino \textbf{Religioso} como uma das cinco áreas de conhecimento do Ensino Fundamental de 09 ( nove ) anos.
Estabelecido como componente curricular de oferta obrigatória nas escolas públicas de Ensino Fundamental, com matrícula facultativa, em diferentes regiões do país, foram elaboradas propostas curriculares, cursos de formação inicial e continuada e materiais didático-pedagógicos que contribuíram para a construção da área do Ensino \textbf{Religioso}, cujas natureza e finalidades pedagógicas são distintas da confessionalidade.
Considerando os marcos normativos e, em conformidade com as competências gerais estabelecidas no âmbito da BNCC, o Ensino \textbf{Religioso} deve atender os seguintes objetivos : 1.
Proporcionar a aprendizagem dos conhecimentos \textbf{religiosos}, culturais e estéticos, a partir das manifestações \textbf{religiosas} percebidas na realidade dos educandos ; 2.
Propiciar conhecimentos sobre o direito à liberdade de consciência e de crença, no constante propósito de promoção dos direitos humanos ; 3.
Desenvolver competências e habilidades que contribuam para o diálogo entre perspectivas \textbf{religiosas} e seculares de vida, exercitando o respeito à liberdade de concepções e o pluralismo de ideias, de acordo com a Constituição Federal ; 4.
Contribuir para que os educandos construam seus sentidos pessoais de vida a partir de valores, princípios éticos e da cidadania.
O conhecimento \textbf{religioso}, objeto da área de Ensino \textbf{Religioso}, é produzido no âmbito das diferentes áreas do conhecimento científico das Ciências Humanas e Sociais, notadamente da(s) Ciência(s) da(s) Religião(ões ).
Essas Ciências investigam a manifestação dos fenômenos \textbf{religiosos} em diferentes culturas e sociedades enquanto um dos bens simbólicos resultantes da busca humana por respostas aos enigmas do mundo, da vida e da morte.
De modo singular, complexo e diverso, esses fenômenos alicerçaram distintos sentidos e significados de vida e diversas ideias de divindade(s), em torno dos quais se organizaram cosmovisões, linguagens, saberes, crenças, mitologias, narrativas, textos, símbolos, ritos, doutrinas, tradições, movimentos, práticas e princípios éticos e morais.
Os fenômenos \textbf{religiosos} em suas múltiplas manifestações são parte integrante do substrato cultural da humanidade.
Cabe ao Ensino \textbf{Religioso} tratar os conhecimentos \textbf{religiosos} a partir de pressupostos éticos e científicos, sem privilégio de nenhuma crença ou convicção.
Isso implica abordar esses conhecimentos com base nas diversas culturas e tradições \textbf{religiosas}, sem desconsiderar a existência de filosofias seculares de vida.
No Ensino Fundamental, o Ensino \textbf{Religioso} adota a pesquisa e o diálogo como princípios mediadores e articuladores dos processos de observação, identificação, análise, apropriação e ressignificação de saberes, visando o desenvolvimento de competências específicas.
Dessa maneira, busca problematizar representações sociais preconceituosas sobre o outro, com o intuito de combater a \textbf{intolerância}, a \textbf{discriminação} e a \textbf{exclusão}.
Por isso, a interculturalidade e a ética da alteridade constituem fundamentos teóricos e pedagógicos do Ensino \textbf{Religioso}, porque favorecem o reconhecimento e respeito às histórias, memórias, crenças, convicções e valores de diferentes culturas, tradições \textbf{religiosas} e filosofias de vida.
O Ensino \textbf{Religioso} busca construir, por meio do estudo dos conhecimentos \textbf{religiosos} e das filosofias de vida, atitudes de reconhecimento e respeito às alteridades.
Trata -se de um espaço de aprendizagens, experiências pedagógicas, intercâmbios e diálogos permanentes, que visam o acolhimento das \textbf{identidades} culturais, \textbf{religiosas} ou não, na perspectiva da interculturalidade, direitos humanos e cultura da paz.
Tais finalidades se articulam aos elementos da formação integral dos estudantes, na medida em que fomentam a aprendizagem da convivência democrática e cidadã, princípio básico à vida em sociedade.
Considerando esses pressupostos, e em articulação com as competências gerais da Educação Básica, a área de Ensino \textbf{Religioso} – e, por consequência, o componente curricular de Ensino \textbf{Religioso} –, devem garantir aos alunos o desenvolvimento de competências específicas.

% matched lemmas: antropológico, discriminação, estado, exclusão, identidade, intolerância, religioso, território
\end{textsample}
